\documentclass[12pt,a4paper]{article}

\usepackage[utf8]{inputenc}
\usepackage[T1]{fontenc}
\usepackage[french]{babel}
\usepackage{amsmath,amssymb,amsthm}
\usepackage{geometry}
\usepackage{booktabs}
\usepackage{longtable}
\usepackage{array}
\usepackage{hyperref}
\usepackage{listings}
\usepackage{graphicx}
\usepackage{xcolor}
\usepackage{enumitem}
\usepackage{pdflscape}
\usepackage{makecell}

\geometry{margin=2.5cm}
\hypersetup{
  colorlinks=true,
  linkcolor=blue,
  urlcolor=blue,
  citecolor=blue,
  pdftitle={Etude de la complexite d'une plateforme de planification basee sur DSL et MiniZinc}
}

\lstset{
  basicstyle=\ttfamily\small,
  frame=single,
  columns=fullflexible,
  breaklines=true
}

\newtheorem{definition}{Définition}[section]
\newtheorem{proposition}{Proposition}[section]
\newtheorem{remark}{Remarque}[section]

\newcommand{\N}{\mathbb{N}}

\title{Étude de la complexité d'une plateforme de planification basée sur DSL et MiniZinc}
\author{LOMOFOUET NDONGMO H. Jauress}
\date{\today}

\begin{document}

\maketitle
\thispagestyle{empty}
\clearpage

\begin{abstract}
Ce rapport étudie la complexité d'une plateforme de planification construite autour d'un DSL et d'une génération vers MiniZinc. L'objectif est d'expliquer, de manière théorique et méthodologique, pourquoi certaines instances se résolvent rapidement alors que d'autres deviennent coûteuses ou conduisent à des \textit{timeouts}. L'analyse distingue explicitement : (i) la complexité logicielle du pipeline (parsing, AST, génération, flattening, appel solveur, restitution), (ii) la complexité combinatoire intrinsèque du problème de planification.

Le document fournit des définitions formelles, des propositions avec preuves, des bornes asymptotiques, une analyse de la taille du modèle généré, une borne naïve de l'espace de recherche, une preuve de NP-difficulté par réduction depuis la coloration de graphe, et une discussion comparative des solveurs (HiGHS, Gecode, Chuffed, OR-Tools CP-SAT). Les résultats expérimentaux n'étant pas fournis, une méthodologie complète de mesure est proposée, avec des tableaux LaTeX prêts à être remplis.
\end{abstract}

\tableofcontents
\clearpage

\section{Introduction}

\subsection{Objectif général du rapport}
L'objectif est d'analyser la complexité d'une plateforme de planification basée sur un DSL, afin de comprendre pourquoi certaines instances sont faciles à résoudre tandis que d'autres sont lentes ou conduisent à un \textit{timeout}. Le temps total peut être décomposé en :
\begin{equation}
T_{\text{total}} = T_{\text{parse}} + T_{\text{AST}} + T_{\text{generation}} + T_{\text{flattening}} + T_{\text{resolution}} + T_{\text{output}}.
\end{equation}

Dans la pratique, la composante dominante est généralement
\begin{equation}
T_{\text{resolution}},
\end{equation}
car la recherche combinatoire est souvent le goulet d'étranglement.

\subsection{Deux niveaux de complexité à distinguer}
\begin{enumerate}[label=\arabic*.]
  \item \textbf{Complexité de la plateforme logicielle :} parsing, construction d'AST, génération MiniZinc, lancement solveur, lecture de la sortie.
  \item \textbf{Complexité du problème combinatoire :} affectations activités-créneaux, activités-ressources, rôles, respect des contraintes, optimisation des préférences.
\end{enumerate}

\subsection{Cadre méthodologique}
Ce rapport distingue explicitement trois types d'énoncés :
\begin{itemize}
  \item \textbf{Résultats théoriques démontrés} (propositions + preuves) ;
  \item \textbf{Estimations asymptotiques} (ordres de grandeur) ;
  \item \textbf{Observations empiriques} (à mesurer expérimentalement, sans résultats inventés).
\end{itemize}

\subsection{Notions cruciales de complexité à retenir}
\begin{itemize}
  \item \textbf{Distinction structurelle :} la génération du modèle est typiquement polynomiale, la résolution peut être exponentielle.
  \item \textbf{Point dominant :} dans la plupart des cas, $T_{\text{resolution}}$ domine le temps total.
  \item \textbf{Explosion combinatoire :} les familles en $O(RA^2)$ (notamment non-chevauchement) et $O(RTA)$ (contraintes par slot) peuvent croître très vite.
  \item \textbf{Espace de recherche :} la borne naïve $T^A \times 2^{AR} \times 2^{AKR}$ explique la difficulté potentielle.
  \item \textbf{Théorie de la difficulté :} sous hypothèses raisonnables, la faisabilité est NP-difficile.
  \item \textbf{Optimisation :} prouver l'optimalité est au moins aussi difficile que trouver une solution faisable.
  \item \textbf{Prédiction :} un score de complexité permet une estimation de risque, pas une garantie de temps exact.
\end{itemize}

\section{Rappels sur la complexité algorithmique}

\subsection{Définitions fondamentales}
\begin{definition}[Complexité temporelle]
La complexité temporelle d'un algorithme est la fonction qui associe à la taille de l'entrée le nombre d'opérations élémentaires nécessaires à son exécution.
\end{definition}

\begin{definition}[Complexité spatiale]
La complexité spatiale d'un algorithme est la quantité de mémoire supplémentaire requise en fonction de la taille de l'entrée.
\end{definition}

\begin{definition}[Notations asymptotiques]
Soient deux fonctions positives $f$ et $g$.
\begin{itemize}
  \item $f(n) \in O(g(n))$ s'il existe $c>0$ et $n_0$ tels que $f(n) \le c g(n)$ pour tout $n\ge n_0$.
  \item $f(n) \in \Omega(g(n))$ s'il existe $c>0$ et $n_0$ tels que $f(n) \ge c g(n)$ pour tout $n\ge n_0$.
  \item $f(n) \in \Theta(g(n))$ si $f(n) \in O(g(n))$ et $f(n) \in \Omega(g(n))$.
\end{itemize}
\end{definition}

\subsection{Pire cas, meilleur cas, cas moyen}
\begin{itemize}
  \item \textbf{Pire cas :} borne supérieure sur toutes les entrées de taille $n$.
  \item \textbf{Meilleur cas :} coût minimal sur les entrées de taille $n$.
  \item \textbf{Cas moyen :} espérance du coût sur une distribution d'entrées.
\end{itemize}

\subsection{Croissances usuelles}
Ordres de grandeur classiques : $O(1)$, $O(n)$, $O(n\log n)$, $O(n^2)$, $O(2^n)$, $O(n!)$. Les classes exponentielles et factorielles deviennent rapidement impraticables.

\begin{longtable}{@{}rrrrr@{}}
\caption{Comparaison qualitative de croissances}\label{tab:croissance}\\
\toprule
$n$ & $O(n)$ & $O(n^2)$ & $O(2^n)$ & $O(n!)$ \\
\midrule
\endfirsthead
\toprule
$n$ & $O(n)$ & $O(n^2)$ & $O(2^n)$ & $O(n!)$ \\
\midrule
\endhead
5  & 5   & 25   & 32     & 120 \\
10 & 10  & 100  & 1024   & 3\,628\,800 \\
20 & 20  & 400  & 1\,048\,576 & $2.43\times 10^{18}$ \\
30 & 30  & 900  & $1.07\times 10^{9}$ & $2.65\times 10^{32}$ \\
\bottomrule
\end{longtable}

\subsection{Complexité théorique vs temps observé}
La complexité asymptotique ne capture ni les constantes cachées, ni l'implémentation, ni l'architecture matérielle, ni les heuristiques du solveur. Elle fournit un cadre de comparaison, pas une prédiction exacte en secondes.

\section{Architecture de la plateforme}

\subsection{Pipeline formel}
Le pipeline considéré est :
\begin{center}
\texttt{DSL $\rightarrow$ Parser $\rightarrow$ AST $\rightarrow$ Générateur MiniZinc $\rightarrow$ Modèle MiniZinc $\rightarrow$ Flattening $\rightarrow$ FlatZinc/Représentation interne $\rightarrow$ Solveur $\rightarrow$ Solution $\rightarrow$ Planning exploitable}
\end{center}

\subsection{Entrées, sorties et coûts par étape}
\begin{longtable}{@{}p{2.4cm}p{2.3cm}p{2.4cm}p{3.2cm}p{3.2cm}p{2.0cm}@{}}
\caption{Analyse des étapes du pipeline}\label{tab:pipeline}\\
\toprule
Étape & Entrée & Sortie & Complexité temporelle (approx.) & Complexité spatiale (approx.) & Remarque \\
\midrule
\endfirsthead
\toprule
Étape & Entrée & Sortie & Complexité temporelle (approx.) & Complexité spatiale (approx.) & Remarque \\
\midrule
\endhead
Parsing & Fichier DSL de taille $L$ & AST & $O(L)$ (ou $O(L\log L)$ selon méthode) & $O(L)$ & Dépend du parseur \\
Construction AST & Tokens/règles & AST enrichi & $O(L)$ & $O(L)$ & Structures syntaxiques \\
Génération MiniZinc & AST & Fichier MZN de taille $M$ & $O(M)$ & $O(M)$ & Écriture textuelle \\
Flattening & MZN & FlatZinc/interne & variable, souvent polynomial mais dépendante du modèle & variable & Pré-traitement solveur \\
Résolution & FlatZinc/interne & solution/infaisabilité/optimalité & potentiellement exponentielle (pire cas) & élevée possible & Étape dominante \\
Sortie & Solution solveur & planning exploitable & $O(|\text{solution}|)$ à $O(M)$ & faible à modérée & Sérialisation \\
\bottomrule
\end{longtable}

\subsection{Remarque sur le flattening}
Le coût du flattening dépend fortement des transformations (décomposition de contraintes globales, expansion d'indices, linéarisation éventuelle). Il est donc préférable de le mesurer empiriquement en plus de l'estimer.

\section{Paramètres de taille d'une instance}

\subsection{Notations}
On introduit les paramètres :
\begin{align*}
A &:= \text{nombre d'activités},\\
R &:= \text{nombre total de ressources},\\
K &:= \text{nombre de rôles},\\
D &:= \text{nombre de jours},\\
S &:= \text{nombre de slots par jour},\\
T &:= D\times S \quad \text{(nombre total de slots)},\\
C &:= \text{nombre total de contraintes explicites DSL},\\
P &:= \text{nombre de préférences / contraintes souples},\\
Q &:= \text{nombre de précédences},\\
W &:= \text{nombre de fenêtres temporelles},\\
E &:= \text{nombre de contraintes d'exclusion/non-chevauchement},\\
G &:= \text{nombre de groupes/types de ressources (si utilisé)},\\
L &:= \text{taille du DSL (caractères ou tokens)}.
\end{align*}

\subsection{Rôle des paramètres}
\begin{itemize}
  \item Paramètres influençant surtout la génération : $L$, $C$, structure de l'AST, schéma de traduction.
  \item Paramètres influençant fortement la résolution : $A$, $R$, $K$, $T$, $P$, $Q$, densité des interactions.
\end{itemize}

Les paramètres critiques sont généralement $A$, $R$, $K$, $T$, $P$ :
\begin{itemize}
  \item $A$ : de nombreuses familles de contraintes croissent en $A^2$ ;
  \item $R$ : multiplication des affectations ;
  \item $K$ : multiplication des variables de rôle ;
  \item $T$ : croissance des choix temporels ;
  \item $P$ : coût supplémentaire de l'optimisation.
\end{itemize}

\section{Taille du modèle MiniZinc généré}

\subsection{Variables principales}
Dans le schéma considéré :
\begin{itemize}
  \item \texttt{start\_time[a]} : une variable de début par activité ;
  \item \texttt{assignment[a,r]} : booléen activité--ressource ;
  \item \texttt{role\_assignment[a,k,r]} : booléen activité--rôle--ressource.
\end{itemize}

\begin{proposition}[Nombre de variables start\_time]
$N_{\text{start}} = A$.
\end{proposition}
\begin{proof}
Par définition du modèle, il existe exactement une variable \texttt{start\_time} pour chaque activité. Comme l'ensemble des activités a cardinal $A$, le nombre total est $A$.
\end{proof}

\begin{proposition}[Nombre de variables assignment]
$N_{\text{assign}} = A\times R$.
\end{proposition}
\begin{proof}
Une variable \texttt{assignment[a,r]} est créée pour chaque couple $(a,r)$ avec $a\in\{1,\dots,A\}$ et $r\in\{1,\dots,R\}$. Le nombre de couples est le produit des cardinaux : $A\times R$.
\end{proof}

\begin{proposition}[Nombre de variables role\_assignment]
$N_{\text{role}} = A\times K\times R$.
\end{proposition}
\begin{proof}
Une variable \texttt{role\_assignment[a,k,r]} est créée pour chaque triplet $(a,k,r)$, avec $a$ activité, $k$ rôle, $r$ ressource. Le nombre total de triplets vaut $A\times K\times R$.
\end{proof}

\begin{proposition}[Nombre total de variables principales]
\begin{equation}
N_{\text{variables}} = A + A\times R + A\times K\times R.
\end{equation}
Asymptotiquement, $N_{\text{variables}}\in O(AKR)$ lorsque $A,R,K$ varient.
\end{proposition}
\begin{proof}
Le total est la somme de trois familles disjointes : \texttt{start\_time}, \texttt{assignment}, \texttt{role\_assignment}. En additionnant les trois cardinalités déjà démontrées :
\[
N_{\text{variables}} = N_{\text{start}} + N_{\text{assign}} + N_{\text{role}} = A + AR + AKR.
\]
Comme $AKR$ domine asymptotiquement lorsque les trois paramètres croissent, la borne $O(AKR)$ suit.
\end{proof}

\subsection{Exemple numérique}
Pour $A=30$, $R=10$, $K=3$ :
\begin{align*}
N_{\text{start}} &= 30,\\
N_{\text{assign}} &= 30\times 10 = 300,\\
N_{\text{role}} &= 30\times 3\times 10 = 900,\\
N_{\text{variables}} &= 30 + 300 + 900 = 1230.
\end{align*}

\section{Analyse de la complexité temporelle}

\subsection{Nombre de contraintes générées par famille}
\begin{longtable}{@{}p{3.0cm}p{3.0cm}p{3.4cm}p{3.2cm}p{2.6cm}@{}}
\caption{Estimations du nombre de contraintes par famille}\label{tab:contraintes-familles}\\
\toprule
Famille & Rôle & Paramètres & Nombre estimé & Asymptotique \\
\midrule
\endfirsthead
\toprule
Famille & Rôle & Paramètres & Nombre estimé & Asymptotique \\
\midrule
\endhead
Domaine activités & Borner dates, domaines & $A,T$ & $\sim A$ à $\sim A\log T$ selon encodage & $O(A)$ à $O(AT)$ \\
Disponibilité ressources & Interdire ressources indisponibles & $A,R,W,T$ & dépend de granularité temporelle & souvent $O(ARW)$ ou $O(ART)$ \\
Affectation activité-ressource & Lier activités et ressources & $A,R$ & $A\times R$ & $O(AR)$ \\
Rôles & Affecter rôle par activité/ressource & $A,K,R$ & $A\times K\times R$ & $O(AKR)$ \\
Cardinalité & Nombre min/max de ressources & $A,R$ & $\sim A$ à $\sim AR$ & $O(A)+O(AR)$ \\
Capacité & Limites cumulées & $A,R,T$ & dépend modélisation & souvent $O(ART)$ \\
Non-chevauchement & Empêcher conflits temporels & $A,R$ & $R\times \binom{A}{2}$ & $O(RA^2)$ \\
Précédence & Ordres partiels & $Q$ ou $A_1,A_2$ & $Q$ ou $A_1A_2$ & $O(Q)$ ou $O(A_1A_2)$ \\
Avant/Après temporel & Fenêtres, délais, lags & $A,W$ & $\sim W$ à $\sim AW$ & $O(W)+O(AW)$ \\
Par slot & Vérifications par créneau & $R,T,A$ & $R\times T\times A$ & $O(RTA)$ \\
Préférences/pénalités & Objectif souple & $P$ & $P$ & $O(P)$ (structure), impact solveur > linéaire possible \\
\bottomrule
\end{longtable}

\subsection{Contraintes de rôle}
\begin{proposition}
Si chaque activité peut mobiliser $K$ rôles et $R$ ressources, alors le nombre de relations de rôle est $A\times K\times R$, donc en $O(AKR)$.
\end{proposition}
\begin{proof}
Chaque relation est indexée par un triplet $(a,k,r)$. Le nombre de triplets possibles est le produit des cardinalités, soit $A\cdot K\cdot R$. La borne asymptotique est $O(AKR)$.
\end{proof}

\subsection{Contraintes d'affectation}
\begin{proposition}
Le nombre de contraintes d'affectation activité--ressource est $A\times R$, donc en $O(AR)$.
\end{proposition}
\begin{proof}
Une contrainte (ou une famille d'atomes linéaires/logiques équivalente) est associée à chaque couple $(a,r)$. Le nombre total de couples est $AR$.
\end{proof}

\subsection{Contraintes de non-chevauchement}
\begin{proposition}
Si le modèle impose une contrainte de non-chevauchement pour chaque paire d'activités et pour chaque ressource, alors le nombre de contraintes de non-chevauchement est
\begin{equation}
N_{\text{overlap}} = R\times \frac{A(A-1)}{2},
\end{equation}
et donc $N_{\text{overlap}}\in O(RA^2)$.
\end{proposition}
\begin{proof}
Le nombre de paires distinctes d'activités vaut $\binom{A}{2}=A(A-1)/2$. Pour chaque ressource, il faut imposer la non-superposition sur chacune de ces paires. En multipliant par $R$ :
\[
N_{\text{overlap}} = R\binom{A}{2}=R\frac{A(A-1)}{2}.
\]
Comme $A(A-1)/2\in\Theta(A^2)$, on obtient $N_{\text{overlap}}\in O(RA^2)$.
\end{proof}

\begin{remark}
Cette famille est critique : si $A$ double, le terme quadratique est multiplié approximativement par $4$ (à $R$ fixé).
\end{remark}

\subsection{Contraintes par slot}
Si une vérification est réalisée pour chaque triplet $(r,t,a)$, la taille est $R\times T\times A$, donc $O(RTA)$. Cette famille devient coûteuse quand $T=D\times S$ croît (horizon long ou granularité fine).

\subsection{Contraintes de précédence}
Si les précédences sont explicitement données sous forme de $Q$ arcs, la taille est $O(Q)$. En revanche, si elles sont générées entre deux ensembles d'activités $A_1$ et $A_2$, on peut atteindre $O(A_1A_2)$.

\subsection{Contraintes de préférences}
Le nombre de contraintes de préférence est souvent $O(P)$. Cependant, leur impact sur $T_{\text{resolution}}$ peut être majeur : le solveur doit explorer et prouver l'optimalité, pas seulement trouver un point faisable.

\subsection{Taille du modèle généré et coût de génération}
On note $M$ la taille (textuelle/structurelle) du modèle MiniZinc généré.

\begin{proposition}[Taille approximative du modèle]
Sous les hypothèses d'encodage usuelles ci-dessus,
\begin{equation}
M = O\big(A + AR + AKR + RA^2 + RTA + Q + P\big).
\end{equation}
\end{proposition}
\begin{proof}
$M$ est la somme des contributions des familles de variables et de contraintes. Les ordres dominants déjà établis sont : $A$, $AR$, $AKR$, $RA^2$, $RTA$, $Q$, $P$. La somme de ces termes donne la borne annoncée.
\end{proof}

\begin{proposition}[Complexité temporelle de la génération]
\begin{equation}
T_{\text{generation}} \in O(M).
\end{equation}
\end{proposition}
\begin{proof}
La génération consiste à parcourir des structures intermédiaires et écrire des fragments textuels de variables/contraintes. Le coût unitaire est proportionnel à la taille des fragments produits. En agrégeant sur l'ensemble des éléments générés, le coût total est proportionnel à la taille totale produite $M$.
\end{proof}

\begin{remark}
Cette étape est en général polynomiale en les paramètres d'entrée, contrairement à la résolution qui peut être exponentielle en pire cas.
\end{remark}

\section{Analyse de la complexité spatiale}

\subsection{Mémoire avant solveur}
Avant résolution, la mémoire inclut : fichier DSL, AST, structures intermédiaires, modèle MiniZinc, et éventuellement buffers de sortie.

\begin{proposition}[Complexité spatiale de génération]
Une estimation est
\begin{equation}
S_{\text{generation}} = O\big(L + A + R + K + T + M\big).
\end{equation}
En pratique, $M$ domine souvent, donc $S_{\text{generation}}\in O(M)$.
\end{proposition}
\begin{proof}
Chaque composant occupe un volume au plus linéaire dans sa taille logique : le DSL en $L$, les métadonnées de taille $A,R,K,T$, puis le modèle produit en $M$. La somme des contributions fournit la borne. Lorsque $M$ est significativement plus grand que les autres termes, il devient dominant.
\end{proof}

\subsection{Discussion pratique}
La mémoire du solveur n'est pas incluse dans $S_{\text{generation}}$ et peut être substantielle (structures de propagation, nœuds de recherche, coupes, clauses apprises selon solveur).

\section{Complexité théorique du problème de planification}

\subsection{Espace de recherche naïf}
\begin{proposition}[Borne naïve de l'espace de recherche]
Si chaque activité choisit un slot parmi $T$, chaque couple activité--ressource est binaire, et chaque triplet activité--rôle--ressource est binaire, alors :
\begin{equation}
E_{\text{search}} = T^A\times 2^{AR}\times 2^{AKR} = T^A\times 2^{AR(K+1)}.
\end{equation}
\end{proposition}
\begin{proof}
Par produit cartésien des choix bruts :
\begin{itemize}
  \item temps : $T$ choix par activité, soit $T^A$ affectations ;
  \item affectations activité--ressource : $AR$ booléens indépendants en borne naïve, soit $2^{AR}$ ;
  \item affectations rôle : $AKR$ booléens, soit $2^{AKR}$.
\end{itemize}
Le nombre total d'états bruts est le produit des cardinalités de ces ensembles de choix, d'où la formule.
\end{proof}

\begin{remark}
Cette borne est pessimiste : la propagation, les relaxations, le \textit{branch-and-bound}, le \textit{branch-and-cut}, l'apprentissage de clauses, les heuristiques et les détections d'infaisabilité réduisent fortement l'espace effectivement exploré.
\end{remark}

\subsection{NP-difficulté par réduction depuis la coloration}
\begin{proposition}[NP-difficulté]
Sous les hypothèses suivantes :
\begin{itemize}
  \item le DSL exprime des contraintes d'incompatibilité/non-simultanéité ;
  \item les slots sont discrets ;
  \item les activités ont une durée compatible avec une occupation d'un slot ;
  \item une solution candidate se vérifie en temps polynomial ;
\end{itemize}
la version décisionnelle du problème de faisabilité du planning est NP-difficile. De plus, elle est NP-complète si la vérification est polynomiale.
\end{proposition}
\begin{proof}
On réduit le problème de $k$-coloration (NP-complet pour $k\ge 3$) au problème de planification.

Soit un graphe $G=(V,E)$ et un entier $k\ge 3$.
\begin{itemize}
  \item Chaque sommet $v\in V$ devient une activité $a_v$.
  \item Les $k$ couleurs deviennent $k$ slots temporels.
  \item Chaque activité doit être placée dans exactement un slot.
  \item Pour chaque arête $(u,v)\in E$, on impose la contrainte : $a_u$ et $a_v$ ne peuvent pas être dans le même slot.
\end{itemize}

Si $G$ est $k$-coloriable, on affecte à chaque activité le slot correspondant à sa couleur : toutes les contraintes sont satisfaites, donc le planning est faisable.

Réciproquement, si un planning faisable existe, l'affectation des slots aux activités définit une couleur par sommet, et les contraintes sur les arêtes garantissent des couleurs différentes aux extrémités : on obtient une $k$-coloration valide.

La transformation est polynomiale en $|V|+|E|$. Donc la faisabilité de planning est au moins aussi difficile que la $k$-coloration : elle est NP-difficile. Si la vérification d'une solution est polynomiale, le problème appartient à NP, donc il est NP-complet.
\end{proof}

\subsection{Optimisation vs satisfaction}
\begin{proposition}
Le problème d'optimisation de planning est au moins aussi difficile que le problème de faisabilité correspondant.
\end{proposition}
\begin{proof}
Toute instance de faisabilité peut être transformée en instance d'optimisation en ajoutant un objectif constant (ou neutre). Un solveur d'optimisation qui retourne une solution optimale permet alors de décider l'existence d'une solution faisable. Donc un algorithme pour l'optimisation résout la faisabilité : l'optimisation est au moins aussi difficile.
\end{proof}

\section{Influence du choix du solveur}

\subsection{Principe général}
Aucun solveur n'est universellement meilleur : les performances dépendent de la représentation MiniZinc obtenue après flattening.

\begin{longtable}{@{}p{2.2cm}p{2.2cm}p{3.5cm}p{3.4cm}p{3.1cm}@{}}
\caption{Comparaison qualitative de solveurs}\label{tab:solveurs}\\
\toprule
Solveur & Type & Modèle favorable & Points forts & Risques / limites \\
\midrule
\endfirsthead
\toprule
Solveur & Type & Modèle favorable & Points forts & Risques / limites \\
\midrule
\endhead
HiGHS & LP/MIP & Formulations linéaires (continues, entières, binaires) & Très bon sur MILP structurés, coupes et bornes efficaces & Moins naturel pour contraintes CP globales \\
Gecode & CP (domaines finis) & Contraintes globales, propagation forte & Bon contrôle de recherche CP, robuste sur CSP structurés & Peut être moins favorable si modèle essentiellement booléen linéarisé massif \\
Chuffed & CP + LCG & Modèles combinatoires avec implications et conflits utiles à apprendre & Apprentissage de clauses, bonne performance sur certains CP difficiles & Sensible à la formulation et au branching \\
OR-Tools CP-SAT & Hybride CP/SAT/MIP & Beaucoup de booléens + linéaire entier + implications & Très performant sur nombreux modèles combinatoires modernes & Comportement dépendant de l'encodage, parfois instable selon structure \\
\bottomrule
\end{longtable}

\subsection{Lecture pratique}
Le choix du solveur doit être guidé par la forme réelle du modèle aplati, pas seulement par le DSL d'origine.

\section{Score de complexité et prédiction du temps}

\subsection{Pourquoi le temps exact est difficile à prédire}
Deux instances ayant mêmes $(A,R,K,T)$ peuvent avoir des temps très différents. Facteurs : densité de contraintes, nombre de solutions, faisabilité vs infaisabilité, coût de preuve d'infaisabilité, présence d'objectif, symétries, qualité d'encodage, stratégie de recherche, solveur, bornes de domaines, contraintes trop faibles ou trop fortes.

\begin{remark}
Différences conceptuelles importantes :
\begin{itemize}
  \item trouver une solution faisable ;
  \item prouver l'absence de solution ;
  \item prouver l'optimalité d'une solution.
\end{itemize}
Ces trois tâches peuvent avoir des coûts radicalement différents.
\end{remark}

\subsection{Exemples conceptuels}
\textbf{Instance potentiellement facile :} ressources abondantes, nombreux slots, peu de contraintes, nombreuses solutions.

\textbf{Instance potentiellement difficile :} ressources rares, fenêtres serrées, incompatibilités nombreuses, préférences à optimiser, solution rare ou absente.

\subsection{Score simple de complexité}
On propose un indicateur structurel :
\begin{equation}
\text{Score}_{\text{simple}} = A + AR + AKR + RA^2 + RTA + Q + P.
\end{equation}
Interprétation des termes : activités, affectations, rôles, non-chevauchement, contraintes par slot, précédences, préférences.

\subsection{Score pondéré}
\begin{equation}
\text{Score} = w_1A + w_2AR + w_3AKR + w_4RA^2 + w_5RTA + w_6Q + w_7P,
\end{equation}
avec $(w_1,\dots,w_7)$ appris/ajustés expérimentalement par solveur et par famille d'instances.

\subsection{Classification de risque (seuils symboliques)}
\begin{longtable}{@{}lll@{}}
\caption{Classification symbolique du risque de lenteur}\label{tab:classification}\\
\toprule
Classe & Condition symbolique & Interprétation \\
\midrule
\endfirsthead
\toprule
Classe & Condition symbolique & Interprétation \\
\midrule
\endhead
Faible & $\text{Score}<\theta_1$ & Résolution probablement rapide \\
Moyen & $\theta_1\le \text{Score}<\theta_2$ & Résolution probablement acceptable \\
Élevé & $\theta_2\le \text{Score}<\theta_3$ & Risque de lenteur \\
Très élevé & $\text{Score}\ge\theta_3$ & Risque de timeout \\
\bottomrule
\end{longtable}

\subsection{Modèle empirique de prédiction}
On peut modéliser : $T_{\text{solve}}\approx f(\text{Score})$, ou utiliser une régression :
\begin{align}
\log(T_{\text{solve}}+1) ={}& \beta_0 + \beta_1A + \beta_2R + \beta_3K + \beta_4T + \beta_5(AR) \nonumber\\
&+ \beta_6(AKR) + \beta_7(RA^2) + \beta_8(RTA) + \beta_9Q + \beta_{10}P.
\end{align}

L'usage de $\log(T_{\text{solve}}+1)$ se justifie par la forte dispersion des temps (ms à minutes), la stabilisation des écarts et une meilleure modélisation des ordres de grandeur. En pratique, prédire une \textbf{classe} (rapide/modéré/lent/timeout probable) est souvent plus robuste qu'une prédiction exacte en secondes.

\section{Limites de l'approche}

\subsection{Limites théoriques}
\begin{itemize}
  \item Les formules proposées donnent des \textbf{ordres de grandeur}, pas des temps exacts.
  \item Les bornes asymptotiques ignorent les constantes, la structure fine des données et les effets d'implémentation.
  \item Les preuves de NP-difficulté portent sur des hypothèses explicites de modélisation (slots discrets, contraintes d'incompatibilité exprimables, etc.).
\end{itemize}

\subsection{Limites pratiques}
\begin{itemize}
  \item Le comportement réel dépend fortement du solveur et du flattening.
  \item Deux instances de même taille peuvent avoir des temps très différents.
  \item Les solveurs utilisent des heuristiques complexes (branching, coupes, propagation, apprentissage), difficiles à modéliser analytiquement.
  \item Les performances doivent être validées expérimentalement sur un corpus représentatif.
\end{itemize}

\section{Conclusion}

Ce rapport établit que :
\begin{enumerate}[label=\arabic*.]
  \item le parsing et la génération sont en général polynomiaux ;
  \item la taille du modèle peut croître rapidement avec $A,R,K,T$ ;
  \item des familles comme le non-chevauchement peuvent atteindre $O(RA^2)$ ;
  \item l'espace de recherche naïf est exponentiel ;
  \item la faisabilité du planning est NP-difficile sous hypothèses raisonnables ;
  \item l'optimisation est au moins aussi difficile que la satisfaction ;
  \item le temps exact est difficile à prédire ;
  \item un score de complexité et un modèle empirique sont utiles pour estimer le risque ;
  \item les mesures expérimentales sont indispensables ;
  \item le choix du solveur dépend fortement de la représentation MiniZinc effective.
\end{enumerate}

La perspective opérationnelle est donc double : maintenir une analyse théorique rigoureuse pour guider la modélisation, et instrumenter la plateforme pour produire des données expérimentales robustes permettant une prédiction fiable des risques de lenteur ou de \textit{timeout}.

\section{Récapitulatif des résultats démontrés}
\begin{enumerate}[label=\arabic*.]
  \item $N_{\text{start}}=A$.
  \item $N_{\text{assign}}=AR$.
  \item $N_{\text{role}}=AKR$.
  \item $N_{\text{variables}}=A+AR+AKR$.
  \item $N_{\text{overlap}}=R\,A(A-1)/2$.
  \item $N_{\text{overlap}}\in O(RA^2)$.
  \item $M=O(A+AR+AKR+RA^2+RTA+Q+P)$.
  \item $T_{\text{generation}}=O(M)$.
  \item $S_{\text{generation}}=O(L+M)$ (souvent dominée par $O(M)$).
  \item $E_{\text{search}}=T^A\times 2^{AR}\times 2^{AKR}$.
  \item NP-difficulté par réduction depuis la $k$-coloration.
  \item L'optimisation est au moins aussi difficile que la satisfaction.
\end{enumerate}

\section{Références}
\begin{thebibliography}{99}

\bibitem{garey-johnson}
M.~R. Garey and D.~S. Johnson,
\textit{Computers and Intractability: A Guide to the Theory of NP-Completeness},
W. H. Freeman, 1979.

\bibitem{apt-cp}
K.~R. Apt,
\textit{Principles of Constraint Programming},
Cambridge University Press, 2003.

\bibitem{rossi-cp}
F. Rossi, P. van Beek, and T. Walsh (eds.),
\textit{Handbook of Constraint Programming},
Elsevier, 2006.

\bibitem{minizinc-book}
K. Marriott, P. Stuckey, M. Garcia de la Banda, and M. Wallace,
\textit{The MiniZinc Handbook},
Springer, 2016.

\bibitem{minizinc-site}
MiniZinc Project,
\textit{MiniZinc Documentation},
\url{https://www.minizinc.org/}.

\bibitem{highs}
HiGHS Optimization,
\textit{HiGHS Documentation},
\url{https://highs.dev/}.

\bibitem{gecode}
Gecode Team,
\textit{Gecode Documentation},
\url{https://www.gecode.org/}.

\bibitem{chuffed}
Chuffed Solver,
\textit{Chuffed Documentation/Repository},
\url{https://github.com/chuffed/chuffed}.

\bibitem{ortools}
Google OR-Tools,
\textit{CP-SAT Solver Guide},
\url{https://developers.google.com/optimization}.

\end{thebibliography}

\end{document}
